% ID: SEL-D0069-OB-E03
% ESTADO: cotejada. Segunda pasada visual y borrador compilado revisados.
% Fuente: Selectividad/sel_2009_matematicas_aplicadas/mat apli CC SS03.pdf
% SHA-256: 8cf5f487f300c573430722e40f741601da1756c5e8c13e541f697d1620373545
% Página: 2; opción B; ejercicio original 3.

\noindent\underline{Parte I}

En un centro escolar, los alumnos de $2^{\mathrm{o}}$ de Bachillerato pueden cursar, como asignaturas optativas, Estadística o Diseño Asistido por Ordenador (DAO). El $70\%$ de los alumnos estudia Estadística y el resto DAO. Además, el $60\%$ de los alumnos que estudia Estadística son mujeres, y de los alumnos que estudian DAO son hombres el $70\%$.
\begin{enumerate}[label=\alph*)]
\item \textbf{(1 punto)} Elegido un alumno al azar, \textquestiondown{}cuál es la probabilidad de que sea hombre?
\item \textbf{(1 punto)} Sabiendo que se ha seleccionado una mujer, \textquestiondown{}cuál es la probabilidad de que estudie Estadística?
\end{enumerate}
\par\noindent\begin{minipage}{\linewidth}
\noindent\underline{Parte II}

En un estudio de mercado del automóvil en una ciudad se ha tomado una muestra aleatoria de 300 turismos, y se ha encontrado que 75 de ellos tienen motor diésel. Para un nivel de confianza del $94\%$:
\begin{enumerate}[label=\alph*)]
\item \textbf{(1.5 puntos)} Determine un intervalo de confianza de la proporción de turismos que tienen motor diésel en esa ciudad.
\item \textbf{(0.5 puntos)} \textquestiondown{}Cuál es el error máximo de la estimación de la proporción?
\end{enumerate}
\end{minipage}\par
